\chapter{Project Synthesis, Related Work, and Conclusion}
\label[chapter]{chap-project-conclusion}

\section{Project Synthesis}

NetFix AI was developed as an engineering decision-support platform with two intentionally independent intelligence workspaces. \textbf{Drive Test Intelligence} converts large LTE field-measurement datasets into prioritized throughput incidents and then explains those incidents through deterministic, KPI-grounded RCA. \textbf{Planning Intelligence} processes a separate network-planning dataset to assign and validate PCI, RSI, Mod3, and Mod4 constraints, visualize the plan, and support deterministic replanning. The platform layer unifies the user experience, reporting, visualization, orchestration, and AI-assisted interaction without treating the two datasets as if they were spatially or operationally linked.

The project therefore addresses three related engineering bottlenecks: the scale and sparsity of telecom data, the time required to isolate the few cases that deserve investigation, and the difficulty of turning complex KPI or planning evidence into a reproducible engineering decision.

\section{End-to-End Outcomes}

\begin{table}[htbp]
\centering
\small
\begin{tabularx}{\textwidth}{>{\bfseries}p{3.0cm} p{4.2cm} X}
\toprule
Workstream & Demonstrated Scale / Result & Final Engineering Output\\
\midrule
Drive-Test Preprocessing & 64,949 raw observations; 1,365 initial columns; LTE-DL modeling population of 32,238 rows & Cleaned and engineered KPI/session dataset suitable for statistical and ML analysis\\
Anomaly Detection & 3,574 final fused rows; 1,429 strict episodes; 1,382 operational incidents & Prioritized, auditable incident handoff with impact, confidence, persistence, and evidence lineage\\
Root Cause Analysis & Deterministic evaluation of 1,382 incident-level aggregates & Cause ID, supporting KPI evidence, evaluation trace, remediation actions, and guarded LLM narrative\\
Planning Intelligence & 2,554 sectors across 725 sites; zero measured hard-rule violations; 53.36 s full planning runtime and 0.42 s independent validation & Reproducible PCI/RSI/Mod3/Mod4 plan with independent rule validation and exportable results\\
Platform Integration & Two independent technical workspaces in one portal & Dashboards, map views, reports, persistent artifacts, guided workflows, and grounded AI interaction\\
\bottomrule
\end{tabularx}
\caption{Summary of the final NetFix AI engineering outputs.}
\label[tab]{project-final-outcomes}
\end{table}

The most important project outcome is not a single model or dashboard. It is the complete evidence chain: raw measurements are reduced to meaningful observations, observations are scored and consolidated into incidents, incidents are explained through explicit telecom rules, and planning decisions are independently validated before being presented to the engineer.

\section{Positioning Against Related Telecom Work}

NetFix AI sits within a broader industry movement toward automated network analytics, assurance, planning, and optimization. Commercial tools such as Infovista TEMS support drive testing and troubleshooting, while Infovista Planet provides RAN planning and optimization \cite{infovista_tems,infovista_planet}. Ericsson Cognitive Software combines AI-supported planning, performance diagnostics, and automated root-cause analysis across the network lifecycle \cite{ericsson_cognitive}. Nokia MantaRay SON similarly applies AI and automation to RAN planning, provisioning, verification, and optimization \cite{nokia_mantaray}.

Academic work also demonstrates the value of data-driven anomaly detection and RCA in cellular networks. Deep RAN models spatial and temporal KPI behavior for large-scale anomaly detection \cite{tanhatalab2019deepran}; Bordeau-Aubert \textit{et al.} separate anomaly detection from anomaly classification in a modular telecom KPI framework \cite{bordeau2023classification}; Moysen \textit{et al.} combine heterogeneous operational data for anomaly identification and performance forecasting \cite{moysen2020bigdata}; and Hasan \textit{et al.} use graph and transformer models for joint anomaly detection and RCA in 5G RANs \cite{hasan2024simba}.

\begin{table}[htbp]
\centering
\scriptsize
\begin{tabularx}{\textwidth}{>{\bfseries}p{3.0cm} p{4.2cm} X}
\toprule
Reference Platform / Work & Primary Focus & Relation to NetFix AI\\
\midrule
Infovista TEMS + Planet & Drive testing, troubleshooting, RF planning, and optimization & Closest commercial comparison to the project's two main engineering directions; NetFix keeps the drive-test and planning datasets independent while exposing both through one lightweight portal\\
Ericsson Cognitive Software & AI-assisted planning, diagnostics, RCA, and optimization & Similar objective of reducing manual network investigation; NetFix emphasizes transparent statistical scoring and an explicit deterministic RCA knowledge base\\
Nokia MantaRay SON & Automated RAN planning, provisioning, verification, and optimization & Comparable automation goal at carrier scale; NetFix remains engineer-controlled and file-driven rather than closed-loop network automation\\
Deep RAN / telecom KPI classification & Learning-based anomaly detection and classification & Related to the ML validation layer; NetFix deliberately keeps interpretable statistical evidence primary and uses session-safe ML as supporting evidence\\
Simba 5G RCA & Learning-based spatio-temporal anomaly detection and root-cause analysis & Demonstrates a more fully learned RCA direction; NetFix instead finalizes the engineering cause deterministically before using an LLM for narration\\
\bottomrule
\end{tabularx}
\caption{Representative commercial and academic work related to NetFix AI.}
\label[tab]{project-related-work}
\end{table}

The comparison is not intended to claim carrier-scale equivalence. Commercial platforms integrate deeply with live operational systems and have broader technology coverage. NetFix AI instead demonstrates a transparent, reproducible project architecture in which statistical evidence, ML predictions, deterministic RCA rules, planning constraints, and human review remain visible to the engineer.

\section{Limitations and Future Directions}

Several extensions would move the platform closer to production-grade telecom operations:

\begin{itemize}
\item \textbf{Operational data integration:} Incorporate OSS counters, alarms, backhaul state, configuration history, and topology data to strengthen outage and congestion diagnoses.
\item \textbf{Dynamic RCA baselines:} Replace selected static engineering gates with cell-specific or time-dependent baselines once sufficient historical data is available.
\item \textbf{Broader field validation:} Repeat the drive-test and planning evaluations on additional regions, operators, bands, mobility conditions, and network configurations.
\item \textbf{Cross-workspace correlation only when justified:} Future integration between Planning Intelligence and Drive Test Intelligence should require verified cell/site identity and spatial alignment rather than assuming the current independent datasets correspond.
\item \textbf{Stronger LLM verification:} Automatically compare every generated diagnostic statement with the structured RCA evidence object before presenting it to an engineer.
\item \textbf{Controlled automation:} Recommended actions can eventually feed approval-based optimization workflows, but automatic parameter changes should remain gated by deterministic validation and engineer authorization.
\end{itemize}

\section{Final Conclusion}

NetFix AI demonstrates an end-to-end approach to telecom engineering intelligence rather than a single isolated AI model. The Drive Test Intelligence workflow starts with a large, sparse measurement export, preserves the telecom evidence required for diagnosis, detects throughput degradation through interpretable statistical families and session-safe ML, and fuses those signals into row, episode, and operational-incident priorities. The RCA layer then converts the prioritized incidents into explicit physical causes through a maintainable three-JSON knowledge base, retains the exact KPI evidence and rule trace behind each decision, and maps confirmed causes to practical remediation guidance.

In parallel, Planning Intelligence addresses a different engineering problem using an independent dataset. Its deterministic planning and validation pipeline assigns and checks PCI, RSI, Mod3, and Mod4 rules at network scale, producing a reproducible plan with zero measured hard-rule violations in the final evaluated dataset. The platform layer brings these independent workflows into one consistent engineering environment without hiding their assumptions or mixing their evidence.

The resulting system is therefore best understood as an \textbf{engineer-centered telecom intelligence platform}: AI and automation reduce the amount of data that must be inspected manually, but the final technical reasoning remains traceable. The project shows that useful telecom intelligence does not require replacing engineering judgment; it can instead organize evidence, rank what matters, explain why a problem was selected, validate proposed network decisions, and give optimization teams a faster path from raw data to an actionable conclusion.
